\scanpage{23}
\mspage{\handwritten{14}}
\uline{Démonstration de} 7.1. Nous prendrons \deleted{pour} $d_o$ \deleted{le plus petit entier} \deleted{\unsure{positif|qui}} tel que l'on ait \handwritten{(entre autres)}
\[
\deleted{\handwritten{d}_o\unsure{\geq|=|-}n+1}\qquad H^i(X,\underline{O}_X(-d))=0\quad\text{pour tout }d\handwritten{\geq}d_o,\ i<n.
\tag*{\handwritten{(*)}}
\]
Désignons par $U$ l'ouvert de $S_d$ formé des $\deleted{t}s\handwritten{\in}S_d$ tels que $\dim Y_s=n-1$, qui est aussi le plus grand ouvert tel que $\underline{Y}\handwritten{\longrightarrow}S_d$ soit plat au dessus de $U$. Comme on a signalé, le choix de $d_o$ implique que pour tout $s\handwritten{\in}U(k')$, les homomorphismes
\[
H^i(X_{k'},\underline{O}_{X_{k'}})\handwritten{\longrightarrow}H^i(Y_s,\underline{O}_{Y_s})
\tag*{\handwritten{(**)}}
\]
sont des isomorphismes pour $i\deleted{\unsure{\geq}}\handwritten{\leq}n-2$, un monomorphisme pour $i=n-1$. Ceci implique, par EGA III 7, que les\deleted{R} $R^if_*(\underline{O}_{\underline{Y}})$ sont plats sur $U$ et y commutent à l'extension de la base, en particulier au passage aux fibres On fera attention que même si $X$ est lisse et géom. connexe, nous serons obligés de façon essentielle de travailler avec des points $s\handwritten{\in}U$ pour lesquels $Y_s$ n\deleted{e}\handwritten{'}est pas irrédu\unsure{\deleted{c}t|\deleted{t}c}\unsure{\deleted{t}i|\deleted{i}t}\unsure{\deleted{i}b|\deleted{b}i}le et a fortiori non lisse. Ce ne sera qu\deleted{e}\handwritten{'} une fois que 7.1 est démontré qu'on pourra en tirer des conclusions \inserted{diophantiennes} sur des point\deleted{\unsure{e}}s femmés dans un ouvert arbitrairement petit de $S_d$ comme dans 7.2, puis les appliquer comme dans 7.3.

Bien entendu, il suffit de prouver 7.1 après extension arbitraire du corps de base $k$, ce qui nous permettrait en particulier de supposer $k$ algébriquement clos; et la conclusion se renforce \inserted{dans le cas général} automatiquement, en remplaçant le groupe de Galois $G$ envisagé par le sous-groupe image du groupe de Galois provenant de la situation géométrique où $k$ est remplacé par $\bar{k}$. Mais pour ppouver 7.1, nous procèderons \inserted{\handwritten{au contraire}} par voie arithmétique, en notant qu'on peut se réduire au cas où $k$ est fini. Cette réduction se fait \handwritten{de} la façon habituelle, et n'offre pas de difficultés. \deleted{Quitte} On peut alors considérer le\deleted{\unsure{f}}\handwritten{s} frobenius ``géométriques'' opérant sur les $H^i(X,\underline{O}_X)$; si $k$ a $q=p^m$ éléments, c\deleted{'}\handwritten{e} \deleted{est}\inserted{\handwritten{sont}} l\deleted{a}\handwritten{es} puissance\handwritten{s} m-ème\handwritten{s} d\deleted{u}\handwritten{es} frobenius
% archive sheet number: 22
